\section{Introducción}

Los centros de atención que reciben solicitudes de servicio (salas de
urgencias, mesas de ayuda técnica, módulos de atención al cliente) enfrentan
un mismo reto operativo: las solicitudes no deben atenderse en el orden en
que llegan, sino según su nivel de urgencia, y el personal necesita poder
consultar o cancelar solicitudes específicas y revisar las atenciones
realizadas sin tener que reorganizar manualmente toda la fila.

Resolver este problema de manera eficiente requiere estructuras de datos que
soporten, simultáneamente: (i) obtener en todo momento la solicitud más
urgente sin tener que ordenar repetidamente toda la colección, (ii) respetar
el orden de llegada entre solicitudes de igual urgencia, (iii) ubicar y
retirar una solicitud puntual por su identificador, y (iv) registrar las
atenciones en el orden en que ocurren para poder consultarlas y, en una
entrega posterior, revertir la más reciente.

Este documento presenta el diseño e implementación inicial de un
\emph{Sistema de Priorización de Atención} (SPA), que resuelve el problema
anterior combinando estructuras lineales estudiadas en la asignatura: colas
y listas enlazadas para el orden de atención, y una pila para el historial
de atenciones.

La contribución de esta primera entrega consiste en: (a) la formulación del
problema y su alcance, (b) la selección y justificación técnica de las
estructuras de datos empleadas, y (c) una implementación funcional que
permite registrar, atender, buscar y cancelar solicitudes, y consultar el
historial de atenciones.

El resto del documento se organiza así: la Sección~\ref{sec:relacionado}
presenta trabajo relacionado; la Sección~\ref{sec:metodologia} describe los
objetivos, el alcance, los requerimientos y el diseño de la solución,
incluyendo las estructuras de datos seleccionadas; la
Sección~\ref{sec:resultados} reporta el estado actual de la implementación;
la Sección~\ref{sec:discusion} discute las limitaciones conocidas; y la
Sección~\ref{sec:conclusiones} resume el estado del proyecto y el trabajo
futuro.
